\section{Setup Object}

A Setup object represents a machining setup within a MIFX package.

A setup describes the configuration in which machining operations are executed.
This may include setup-level resources such as part, stock, fixture, setup offsets, notes, and setup assemblies.

Each setup is stored as an independent JSON object within the \texttt{setups} directory.

Example location:

\begin{verbatim}
setups/set-1/setup.json
\end{verbatim}

\subsection{Setup Structure}

A Setup object typically contains:

\begin{itemize}
\item a unique identifier
\item a setup index
\item a descriptive name
\item optional artifacts describing setup resources
\item an ordered list of operations executed within the setup
\end{itemize}

Example:

\begin{verbatim}
{
  "id": "set-1",
  "index": 1,
  "name": "1st setup",

  "artifacts": [
    {
      "role": "notes",
      "kind": "json",
      "path": "notes.json",
      "present": true
    },
    {
      "role": "setup_offset",
      "kind": "json",
      "path": "offset.json",
      "present": true
    }
  ],

  "operations": [
    {
      "id": "op-1",
      "role": "toolpath",
      "kind": "mill",
      "path": "operations/op-1.json"
    },
    {
      "id": "op-2",
      "role": "toolpath",
      "kind": "mill",
      "path": "operations/op-2.json"
    }
  ]
}
\end{verbatim}

\subsection{Setup Identifier}

Each setup must define an identifier using the \texttt{id} field.

Example:

\begin{verbatim}
"id": "set-1"
\end{verbatim}

\subsection{Setup Index}

The \texttt{index} field indicates the logical order of the setup within the machining job.

Example:

\begin{verbatim}
"index": 1
\end{verbatim}

This value is typically determined by the exporting CAM system.

\subsection{Setup Artifacts}

The Setup object may reference additional setup-level resources using artifacts.

Common setup artifacts include:

\begin{itemize}
\item notes
\item setup offsets
\item part definitions
\item stock definitions
\item fixture definitions
\item setup assemblies
\end{itemize}

These resources are stored as separate files within the setup directory.

\subsection{Operations}

The \texttt{operations} array defines the machining operations executed in the setup.

Each entry contains:

\begin{itemize}
\item the operation identifier
\item the operation role
\item the operation kind
\item the relative path to the Operation object
\end{itemize}

Example:

\begin{verbatim}
{
  "id": "op-2",
  "role": "toolpath",
  "kind": "mill",
  "path": "operations/op-2.json"
}
\end{verbatim}

The order of entries in the \texttt{operations} array represents the execution order of the machining process.
